\documentclass[a4paper,10pt, twoside]{article}
\usepackage[margin=1.4in]{geometry}
\usepackage[swedish]{babel}
\usepackage[utf8]{inputenc}
\usepackage{titlesec}
\usepackage{titling}
\usepackage{todonotes}



\setlength{\parskip}{1em}
\setlength{\parindent}{0pt}
\titlespacing{\section}{0pt}{\parskip}{-\parskip}
\titlespacing{\subsection}{0pt}{\parskip}{-\parskip}
\titlespacing{\subsubsection}{0pt}{\parskip}{-\parskip}
\titlespacing{\part}{0pt}{\parskip}{-\parskip}

\usepackage{multicol}
\usepackage[backend=bibtex,sorting=none]{biblatex}

\def\ftitle{Testplan}
\def\fversion{1.2}

\addbibresource{bib-testplan.bib}

\externaldocument[krav-]{../Kravspec/kravspec}
\externaldocument[arki-]{../Arkitekturdokument/arkitekturdokument}

\begin{document}
\begin{titlepage} % Suppresses displaying the page number on the title page and the subsequent page counts as page 1
	\newcommand{\HRule}{\rule{\linewidth}{0.5mm}} % Defines a new command for horizontal lines, change thickness here

	\center % Centre everything on the page

	%------------------------------------------------
	%	Headings
	%------------------------------------------------

	\textsc{\LARGE Linköpings universitet \\ \vspace{0.2em} Institutionen för datavetenskap }\\[2cm]

    \large\today

    \vspace{1cm}


	%------------------------------------------------
	%	Title
	%------------------------------------------------

	\HRule\\[0.4cm]

	{\huge\bfseries Schemaläggningsstöd för  kirurgi \vspace{.1em} \\ \ftitle }\\[0.4cm] % Title of your document

	\HRule\\[1cm]

	%------------------------------------------------
	%	Author(s)
	%------------------------------------------------

	\begin{minipage}{0.7\textwidth}
			\large
            \emph{Version: \fversion}
            \vspace{1em}

            \textbf{\\Adam Andersson, Niclas Byrsten, \\Björn Hvass, Henrik Lindström, \\Martin Persson, Christoffer Sjöbergsson, \\Tor Utterborn}


            \vspace{1em}

            Handledare: Jonas Wallgren

            Examinator: Kristian Sandahl
	\end{minipage}
	~

	%------------------------------------------------
	%	Logo
	%------------------------------------------------

	%\vfill\vfill
	%\includegraphics[width=0.2\textwidth]{../Templates/Aeon}\\[1cm] % Include a department/university logo - this will require the graphicx package

	%----------------------------------------------------------------------------------------

	\vfill % Push the date up 1/4 of the remaining page

\end{titlepage}


\section*{\begin{center}Projektidentitet\end{center}}
    \vspace{-2.5em}
    \begin{center}
        \begin{tabular}{|c c c|}
        \hline
        Namn & Roll & E-post \\
        \hline
        Adam Andersson& Teamleader & adam.e.a.andersson@gmail.com\\
        \hline
        Niclas Byrsten & Testansvarig & nicby889@student.liu.se\\
        \hline
        Björn Hvass & Konfigurationsansvarig & bjorn.hvass@gmail.com\\
        \hline
        Henrik Lindström & Utvecklingsledare & henli070@student.liu.se\\
        \hline
        Martin Persson & Arkitekturansvarig & marpe902@liu.student.se\\
        \hline
        Christoffer Sjöbersson & Analysansvarig & chrsj812@liu.student.se\\
        \hline
        Tor Utterborn & Dokument- \& Kvalitetsansvarig & tor.utterborn@gmail.com\\
        \hline
        \end{tabular}
    \end{center}

    \begin{center}
        \small
        \textbf{Kund}\\Region Östergötland, 581 91 Linköping.

        \textbf{Kontaktperson hos kund}\\
        Gunnar Nordqvist, IT-arkitekt, 010-1030698, Gunnar.Nordqvist@regionostergotland.se\\
        Erik Sundvall, Informationsarkitekt, 010-1036252, Erik.Sundvall@regionostergotland.se
    \end{center}

\vspace{9em}



\section*{\begin{center}Dokumentationshistorik\end{center}}
	\begin{center}
        \begin{tabular}{|c c c c |}
        \hline
        Datum & händelse & iteration & version\\
        \hline
        2018-02-14 & Dokument skapas & 1 & 0.1\\
        \hline
        2018-02-19 & Konvertering till \textbf{\LaTeX} & 1 & 0.2\\
        \hline
        2018-03-05 & Version 1.0 färdigställd & 1 & 1.0\\
        \hline
        2018-04-17 & Revision av version 1.0 & 2 & 1.1\\
        \hline
        2018-05-20 & Revision inför slutinlämning & 4 & 1.2\\
        \hline
        \end{tabular}
    \end{center}
\clearpage
\tableofcontents
\clearpage

\section{Introduktion}
Det här dokumentet ska översiktligt beskriva de tester som kommer att genomföras under projektets gång.

\subsection{Mål och syfte}
Testerna är till för att säkerställa att det som levereras uppfyller de krav som finns i kravspecifikation sektion 
\ref{krav-sec:Specifika krav}.

\section{Moduler}
Systemet är uppdelat i två delar, webbservern och webbklienten. Dessa delarna är i sin tur uppdelade i ett antal olika 
moduler. En övergripande systemskiss finns i arkitekturdokumentet under rubrik \ref{arki-sec:Systemskiss}. Mer information
angående webbklienten och dess moduler se sektion \ref{arki-sec:webbklient} medan information angående webbservern finns 
under sektion \ref{arki-sec:webbserver}.

\section{Funktioner}
Denna sektion beskriver de funktioner som ska genomgå testning. Varje funktions användningsområde presenteras tillsammans 
med den indata som funktionen ska behandla. Varje funktion har även blivit tilldelad en integritets nivå som
beskriver hur kritisk funktionen är för systemet. Se kravspecifikationen \ref{krav-subsec:Funktionellakrav} för mer 
ingående information om kraven som berör funktionen.

\subsection{Boka Operation}
\label{sec:Boka Operation}
I denna funktion skall man kunna välja tid, se vilka lediga bokningsbara tider som existerar samt bli informerad om när 
kompetens, material och utrustning saknas. Förberedelser och efterarbete beskrivs med hjälp av KVÅ-kodning.

Beslutsdata
\begin{multicols}{3}
\begin{itemize}
	\item Personnummer
	\item Efternamn, Förnamn
	\item Diagnoskod
	\item Åtgärdskod enligt KVÅ-standard
	\item vilken KVÅ-kod som är huvudåtgärd
	\item Operation av parigt organ (lateralitet)
	\item Datum och tid för beslut om kirurgisk åtgärd
	\item Namn på beslutsfattare inkl. titel
	\item Beslutande klinik
	\item Medicinsk brådskandegrad
	\item Behandlingsnummer
	\item Akut eller elektiv operation
\end{itemize}
\end{multicols}

Operationstider (KVÅ-koder)
\begin{multicols}{3}
\begin{itemize}
	\item Förberedelsetid
	\item Operationstid
	\item Avvecklingstid.
\end{itemize}
\end{multicols}

\subsubsection{Boka sal}
\label{sec: Boka sal}
Det ska gå att välja namn på den lokal man vill boka och välja vilken typ av lokal det är. Man ska även kunna se 
tillgängliga tider för salen. Vid bokning skall man få information över vilka operationer som går att genomföra i lokalen 
samt en ingreppbeskrivning.

Ingående data
\begin{multicols}{2}
\begin{itemize}
	\item Namn på lokal
	\item Typ av lokal
	\item Tillgängliga tider för sal
	\item Möjliga operationer att genomföra i lokalen
	\item Ingreppsbeskrivning
\end{itemize}
\end{multicols}

\subsubsection{Avboka sal}
\label{sec: Avboka sal}
Det ska vara möjligt att avboka en tid som man sedan tidigare bokat.

- Välja tid

\subsection{Inloggning}
\label{sec: Inloggning}
I denna funktion ska följande uppfyllas:
\begin{enumerate}
	\item Det ska krävas användarnamn och lösenord för att komma in i systemet.
  	\item En användare ska ha ett användarnamn och ett lösenord.
  	\item Det ska finnas en databas som kan hantera all ingående data.
\end{enumerate}

\subsection{Tidslinje för operation}
\label{sec: Tidlinje för operation}
Ska presentera den beslutsdata som finns att tillgå.

-	Visa beslutsdata

\subsection{Översikt av beslutade operationer}
\label{sec: Översikt av beslutade operationer}
Ska presentera beslutade operationer.

-	Visa beslutade operationer
\subsection{Schemaöversikt för planerade operationer}
\label{sec: Schemaöversikt för planerade operationer}
En överblick för de operationer som ej blivit beslutade men som är planerade, skall gå att filtrera efter önskemål angivna 
i kravspec.

- Filterfunktion för schemaöversikt

\clearpage
\subsection{Reservera material, lokal, kompetens och utrustning}
\label{sec: Resevera material, lokal, kompetens och utrustning}
Databasfunktion som kan reservera de olika resurser som behövs för operation samt post- eller pre-operativa processer.
\subsubsection{Material}
\label{Material}
Ingående data:
\begin{multicols}{2}
\begin{itemize}
	\item Material
	\item Artikelnamn
	\item Typ av material
	\item Reserverade tider för material som måste steriliseras \\Tidsåtgång: \\ Operationstid + steriliseringstid
\end{itemize}
\end{multicols}

\subsubsection{Utrustning}
\label{sec: Utrustning}
Ingående data
\begin{multicols}{2}
\begin{itemize}
	\item Namn
	\item MTÖ-nummer
	\item Typ av utrustning
	\item Reserverade tider.
\end{itemize}
\end{multicols}

\subsubsection{Kompetenser}
\label{sec: Kompetenser}
Ingående data
\begin{multicols}{2}
\begin{itemize}
	\item Titel
	\item Efternamn, förnamn
	\item Kliniktillhörighet
	\item Klinisk specialitet
	\item Kompetens (Kvå-kod, huvudåtgärd)
	\item Schema
\end{itemize}
\end{multicols}

\section{Tester}
Den här sektionen redogör för de olika testningsprocedurerna som ska utföras för de olika testnivåerna.

\subsection{Enhetstest}
\label{sec:Enhetstest}
\textbf{Syfte:} Efter det att en funktion har implementerats ska koden genomgå ett enhetstest. Testet ska utföras för att 
säkerställa att funktionen fungerar enligt specifikationen \cite{kravspec}. Vidare används testet för att fastställa att 
funktionen är följer designkraven. Testerna kommer att vara av typen black-box.

\textbf{Ingående data:}
\begin{enumerate}
    \item Enheten som ska testas.
    \item Testspecifikation-mall. \emph{(Se appendix \ref{sec: Testspecifikationmall} Testspecifikationmall)}
    \item Testprotokoll-mall. \emph{(Se appendix \ref{sec: Testprotokollmall} Testprotokollmall)}
    \item Bugglista: Den här är listan ska innehålla alla kända problem som systemet har vid testets start, listan kommer 
    att finnas under \emph{issues} på projektets gitlab. Det är sen upp till valfri person att lösa dem.
\end{enumerate}

\textbf{Utförande:} Testet ska utföras av berörda parter samt en oberoende person. Personerna som utför testet ska först 
skapa testspecifikation utifrån beskrivningen i testspecifikationsmallen. När det här är klart så utför personerna testet 
enligt specifikationen \cite{kravspec}. Ett protokoll ska föras enligt mallen för testprotokoll, (se appendix \ref{sec: 
Testprotokollmall} testprotokollmall). Om några problem uppstod som inte kunde lösas på plats ska dessa föras in under 
\emph{issues} i projektets gitlab.

\textbf{Resultat:}
\begin{enumerate}
	\item \label{itm:Testspecifikation-e} Testspecifikation-testtyp-komponent/modul (komponent/modul är namnet på 		   
	komponenten eller modulen som specifikation är skriven för).
    \item \label{itm:Testprotokoll-e} Testprotokoll-testtyp-komponent/modul (komponent/modul är namnet på komponenten 
    eller modulen som specifikation är skriven för).
    
    \item Bugglista: Listan ska vid behov uppdateras.
\end{enumerate}

\textbf{Testkriterium:} Om resultatet från testprotokollet, se punkt \ref{itm:Testprotokoll-e} under \emph{resultat} i 	   
sektion \ref{sec:Enhetstest} ovan, stämmer överens med det förväntade resultatet från testspecifikationen, punkt 
\ref{itm:Testspecifikation-e} under \emph{resultat} i sektion \ref{sec:Enhetstest} ovan, så anses komponenten klarat 
testet.

\textbf{Enhetstester under iteration 4 :} 
\\ Under iteration 4 kommer följande enhetstester att genomföras. Tidpunkten kommer att bestämmas under nästa 
sprint planering.

\textbf{TestID:} \ref{sec:Boka Operation}
\\ \textbf{Tidpunkt:}
\\ \textbf{Syfte och funktion:} Säkerhetsställa att det går att välja tid, se lediga bokningsbara tider, bli informerad 
när kompetens, material eller utrustning saknas.
  
\textbf{TestID:} \ref{sec: Boka sal} 
\\ \textbf{Tidpunkt:}
\\ \textbf{Syfte och funktion:} Namn på lokal och typ av lokal ska kunna bokas. Det ska även gå att se tillgängliga tider 
för lokalen.

\textbf{TestID:} \ref{sec: Avboka sal}
\\ \textbf{Tidpunkt:}
\\ \textbf{Syfte och funktion:} Ska vara möjligt att avboka en bokad lokal.
  
\textbf{TestID:} \ref{sec: Inloggning}
\\ \textbf{Tidpunkt:}  
\\ \textbf{Syfte och funktion:} Det ska gå att logga in i systemet med ett giltigt lösenord.
  
\textbf{TestID:} \ref{sec: Tidlinje för operation}
\\ \textbf{Tidpunkt:}
\\ \textbf{Syfte och funktion:} Visa tillgänglig beslutsdata

\textbf{TestID:} \ref{sec: Översikt av beslutade operationer}
\\ \textbf{Tidpunkt:}  
\\ \textbf{Syfte och funktion:} Presentera belutade operationer.

\textbf{TestID:} \ref{sec: Schemaöversikt för planerade operationer}
\\ \textbf{Tidpunkt:}  
\\ \textbf{Syfte och funktion:} Visa planerade men ej beslutade operationer som ska kunna filtreras.
  
\textbf{TestID:} \ref{sec: Resevera material, lokal, kompetens och utrustning} 
\\ \textbf{Tidpunkt:}
\\ \textbf{Syfte och funktion:} Reservera resurser som behövs för operation samt post- eller pre-operativa processer.

\subsection{Integrationstest}
\label{sec:Integrationstest}
\textbf{Syfte:} Ett integrationstest ska utföras för att säkerställa att en modul och dess komponenter fungerar enligt 
specifikationen \cite{kravspec}. Testet säkerställer också att modulen följer designkraven. När funktioner är klara och 
testade och ska sammanfogas ska de integrationstestas.

\textbf{Ingående data:}
\begin{enumerate}
	\item Modulen som ska integrationstestas.
    \item Testprotokoll för alla enheter som finns i modulen som ska integrationstestas.
    \item Testspecifikation-mall. \emph{(Se appendix \ref{sec: Testspecifikationmall} {Testspecifikationmall})}
    \item Testprotokoll-mall. \emph{(Se appendix \ref{sec: Testprotokollmall} Testprotokollmall)}
    \item Bugglista: Den här är listan ska innehålla alla kända problem som systemet har vid testets start, listan kommer 
    att finnas under \emph{issues} på projektets gitlab. Det är sen upp till valfri person att lösa dem.
\end{enumerate}

\textbf{Utförande:}
Testet ska utföras av berörda parter, det bör max finnas en representant per komponent som ska integrationstestas. 
Personerna som utför testet ska börja med att skapa testspecifikation utifrån den beskrivning som finns i 
testspecifikationsmallen. När det är klart så utför personerna testet enligt specifikationen \cite{kravspec}. Ett 
protokoll ska föras enligt mallen för testprotokoll, se appendix \ref{sec: Testprotokollmall} Testprotokollmall. Om några 
problem uppstod som inte kunde lösas på plats ska dessa föras in under \emph{issues} i projektets gitlab.

\textbf{Resultat:}
\begin{enumerate}
\item \label{itm:Testspecifikation-i} Testspecifikation-testtyp-komponent/modul (komponent/modul är namnet på komponenten 
eller modulen som specifikation är skriven för).
\item \label{itm:Testprotokoll-i} Testprotokoll-testtyp-komponent/modul(komponent/modul är namnet på komponenten eller 
modulen som specifikation är skriven för).
\item Bugglista: Listan ska vid behov uppdateras.
\end{enumerate}

\textbf{Testkriterium:}
Om resultatet från testprotokollet, se punkt \ref{itm:Testprotokoll-i} under \emph{resultat} i sektion
\ref{sec:Integrationstest} ovan, stämmer överens med det förväntade resultatet från testspecifikationen, punkt 
\ref{itm:Testspecifikation-i} under \emph{resultat} i sektion \ref{sec:Integrationstest} ovan, så anses komponenten klarat 
testet.

\textbf{Integrationstester under iteration 4}
De integrationstester som kommer att genomföras under iteration 4. Tidpunkt tillkommer. Mer information om testet kommer 
att finnas i det testprotokoll som kommer att skrivas i samband med testet. 

\textbf{TestID:} I.1
\\ \textbf{Tidpunkt:}
\\ \textbf{Syfte och funktion:} Säkerhetsställa att funktioner fungerar då de är sammanfogade.  


\subsection{Systemtest}
\label{sec:Systemtest}
\textbf{Syfte:} Testet ska utföras för att säkerställa att systemet behandlar och utan problem kan utföra alla de 
funktioner som krävs av systemet enligt specifikationen \cite{kravspec}.

\textbf{Ingående data:}
\begin{enumerate}
	\item Systemet som ska testas
    \item Testprotokoll för alla systemets moduler
    \item Testspecifikation-mall. \emph{(Se appendix \ref{sec: Testspecifikationmall} Testspecifikationmall)}
    \item Testprotokoll-mall. \emph{(Se appendix \ref{sec: Testprotokollmall} Testprotokollmall)}
    \item Bugglista: Den här är listan ska innehålla alla kända problem som systemet har vid testets start, listan kommer 
    att finnas under \emph{issues} på projektets gitlab. Det är sen upp till valfri person att lösa dem.
\end{enumerate}

\textbf{Utförande:} Testet ska utföras av hela projektgruppen. Gruppen utser två till tre personer som skapar en 
testspecifikation utifrån beskrivningen i testspecifikationsmallen. När det är klart så utför personerna testet enligt 
specifikationen \cite{kravspec}. Ett protokoll ska föras enligt mallen för testprotokoll, se appendix \ref{sec: 
Testprotokollmall} Testprotokollmall. Om några problem uppstod som inte kunde lösas på plats ska dessa föras in under 
\emph{issues} i projektets gitlab.

\textbf{Resultat:}
\begin{enumerate}
	\item \label{itm:Testspecifikation-s} Testspecifikation-testtyp-komponent/modul (komponent/modul är namnet på 
	komponenten eller modulen som specifikation är skriven för)
    \item \label{itm:Testprotokoll-s} Testprotokoll-testtyp-komponent/modul(komponent/modul är namnet på komponenten 	
    eller modulen som specifikation är skriven för)
    \item Bugglista: Listan ska vid behov uppdateras
\end{enumerate}

\textbf{Testkriterium:} Om resultatet från testprotokollet, se punkt \ref{itm:Testprotokoll-s} under \emph{resultat} i 
sektion \ref{sec:Systemtest} ovan, stämmer överens med det förväntade resultatet från testspecifikationen, punkt 
\ref{itm:Testspecifikation-s} under \emph{resultat} i sektion \ref{sec:Systemtest} ovan, så anses komponenten klarat 
testet.

\textbf{Systemtester under iteration 4}
\\ Under iteration 4 ska ett systemtest på hela systemet genomföras. Tidpunkt tillkommer under nästa sprint planning.

\textbf{TestID:} S.1
\\ \textbf{Tidpunkt:}
\\ \textbf{Syfte och funktion:}

\subsection{Acceptanstest}
\textbf{Syfte:} Två stycken acceptanstest ska genomföras. Ett test ska utföras på systemets grafiska gränssnitt.  Det 
andra testet ska utföras på hela systemet efter att systemtester blivit godkänt.

\textbf{Ingående data:}
\begin{enumerate}
	\item Testprotokoll-system
    \item Modulen(a) som ska testas
    \item Testspecifikation-mall \emph{(Se appendix \ref{sec: Testspecifikationmall} Testspecifikationmall)}
    \item Testprotokoll-mall \emph{(appendix \ref{sec: Testprotokollmall} Testprotokollmall)}
    \item Bugglista: Den här är listan ska innehålla alla kända problem som systemet har vid testets start, listan kommer 
    att finnas under \emph{issues} på projektets gitlab. Det är sen upp till valfri person att lösa dem.
\end{enumerate}

\textbf{Utförande:} Testet ska utföras av hela projektgruppen. Gruppen utser två till tre personer som skapar en 
testspecifikation utifrån beskrivningen i testspecifikationsmallen. När det är klart så utför personerna testet enligt 
specifikationen \cite{kravspec}. Ett protokoll ska föras enligt mallen för testprotokoll, se testprotokoll-mall 
dokumentet. Om några problem uppstod som inte kunde lösas på plats ska dessa föras in under \emph{issues} i projektets
gitlab.

\textbf{Resultat:}
\begin{enumerate}
	\item \label{itm:Testspecifikation-a} Testspecifikation-testtyp-komponent/modul (komponent/modul är namnet på 			
	 		  komponenten eller modulen som specifikation är skriven för)
    \item \label{itm:Testprotokoll-a} Testprotokoll-testtyp-komponent/modul(komponent/modul är namnet på komponenten eller 
    		  modulen som specifikation är skriven för)
    \item Bugglista: Lisan ska vid behov uppdateras
\end{enumerate}

\textbf{Testkriterium:} Testet anses vara godkänt om och endast om alla krav i specifikationen är behandlade och 
uppfyllda.
\clearpage

\textbf{Acceptanstester under iteration 4}
\\ Under iteration 4 ska två stycken acceptanstester genomföras. Tidpunkt tillkommer. 

	\textbf{TestID:} A.1
 \\ \textbf{Tidpunkt:}
 \\ \textbf{Syfte och funktion:} 

 	\textbf{TestID:} A.2
 \\ \textbf{Tidpunkt:}
 \\ \textbf{Syfte och funktion:}

\printbibliography

\newpage

\appendix
\section{Testspecifikationmall}
\label{sec: Testspecifikationmall}

\textbf{TestID}
\\ \textbf{Typ av test:} \emph{[Vilket typ av test är det]}
\\ \textbf{Genomförande:} \emph{[kort beskrivning av sektion]}
\\ \textbf{Testscenarion:} \emph{[Vad är det som ska testas, kan vara flera saker]}
\\ \textbf{Första saken att testa:}
\\ \textbf{Design av komponent:} \emph{[Beskriv design av komponent och om den är uppfylld enligt Kravspec]}
\\ \textbf{Förväntat resultat:} \emph{[Det förväntade resultatet enligt kravspec eller beskrivning av enheten]}

\newpage

\section{Testprotokollmall}
\label{sec: Testprotokollmall}

\textbf{TestID:}
\\ \textbf{Testansvarig:\\Observatör:}
\\ \textbf{Datum:}
\\ \textbf{Resurser:} \emph{[Vilka resurser kommer att användas i testet, med Referenser om så behövs]}
\\ \textbf{Test log:} \emph{[En detaljerad log över vad som skedde under testets gång]\\ T.ex. \\ 14:23 Testet startades \
\ 14:39 Fann en bugg huxflux \\ 15:40 Åtgärdade buggen huxflux \\ 16:00 Testet slutfördes}\
\\ \textbf{Resultat:} \emph{[Vad blev resultat och blev den godkänd]}
\\ \textbf{Buggar:} \emph{[Här så ska alla olösta buggar som hittades under testet listas OBS! de ska också skrivas under 
issues i gitlab]}
\end{document}
